\chapter{תכנון הפתרון וארכיטקטורה}\label{ch:design}

הפתרון בנוי משני חלקים שעובדים יחד: פלטפורמה טכנולוגית שמייצרת ראיות, ותקן אתי שקובע מה עושים עם הראיות. הפרק מתאר את שניהם, מהתמונה הכללית ועד רמת הרכיב.

\section{הקשר המערכת}
\begin{figure}[!htb]
\begin{ltrdiagram}
% context: TikZ / pgfplots source, built inside ltrdiagram (LTR)
\begin{tikzpicture}[font=\sffamily\small,
  core/.style={draw=acc2,fill=soft2,very thick,rounded corners=6pt,minimum width=3.6cm,minimum height=1.6cm,align=center},
  ext/.style={draw=acc,fill=soft,rounded corners=4pt,minimum width=3.0cm,minimum height=0.95cm,align=center},
  ar/.style={{Stealth[length=5pt]}-{Stealth[length=5pt]},thick,ink!70}]
\node[core] (v) at (0,0) {\en{\textbf{VeriLens}}\\\he{פלטפורמת אימות}};
\node[ext] (cms) at (-5,1.6) {\he{מערכת ניהול תוכן}\\\en{(CMS)}};
\node[ext] (app) at (-5,-1.6) {\he{אפליקציית הצלם}\\\he{חתימת \en{C2PA}}};
\node[ext] (rev) at (5,1.6) {\he{חיפוש תמונה הפוך}};
\node[ext] (crl) at (5,0) {\he{רשימות ביטול}\\\en{C2PA}};
\node[ext] (ads) at (5,-1.6) {\he{רשתות פרסום}\\\he{ציון שלמות}};
\node[ext] (read) at (0,-2.9) {\he{קוראים: תווית}\\\he{ותעודת מקור}};
\draw[ar] (cms)--(v) node[midway,above,sloped,font=\sffamily\scriptsize]{\he{תמונה / החלטה}};
\draw[ar] (app)--(v) node[midway,below,sloped,font=\sffamily\scriptsize]{\he{תמונה חתומה}};
\draw[ar] (rev)--(v);
\draw[ar] (crl)--(v);
\draw[ar] (ads)--(v) node[midway,below,sloped,font=\sffamily\scriptsize]{\he{ציון מצרפי}};
\draw[-{Stealth[length=5pt]},thick,ink!70] (v)--(read);
\end{tikzpicture}

\end{ltrdiagram}
\caption{המערכת יושבת בין מערכת ניהול התוכן לבין הקורא, ומתחברת לחמישה גורמים חיצוניים}\label{fig:context}
\end{figure}

לכל אחד מחמשת החיבורים שבאיור יש שורה בטבלה~\ref{tab:icd}, שמגדירה את הממשק.


\begin{table}[!htb]\centering
\caption{מסמך ממשקים מקוצר: חמשת הממשקים החיצוניים}\label{tab:icd}
\small
\begin{tabularx}{\linewidth}{@{}l Y l l@{}}
\toprule
\textbf{ממשק} & \textbf{מה עובר בו} & \textbf{פורמט} & \textbf{זמן תגובה} \\
\midrule
\req{IF-01} מערכת ניהול התוכן & תמונה נכנסת, החלטה ותווית יוצאות & \en{REST}, \en{JSON} & עד \num{60} שנ׳ \\
\req{IF-02} אפליקציית הצלם & תמונה חתומה בתעודת מקור & \en{C2PA} & מיידי \\
\req{IF-03} חיפוש תמונה הפוך & טביעת תמונה יוצאת, הופעות קודמות נכנסות & \en{REST} & עד \num{20} שנ׳ \\
\req{IF-04} רשימות ביטול & סטטוס תעודות & \en{OCSP}, \en{CRL} & עד \num{2} שנ׳ \\
\req{IF-05} רשתות פרסום & ציון שלמות מצרפי לאתר, בלי פרטי תמונות & \en{REST} & עד \num{1} שנ׳ \\
\bottomrule
\end{tabularx}
\end{table}

\section{ארכיטקטורה ברמת מערכת}
איור~\ref{fig:arch} מציג את שש תתי המערכות. העיקרון המרכזי הוא הפרדה בין ייצור ראיות (שכבות \num{1} עד \num{3}) לבין החלטה (שכבה \num{4}): כל שכבת ראיות מוסיפה ממצאים לתיק התמונה, ורק מנוע השקלול ממיר אותם לציון ולהמלצה.

\begin{figure}[!htb]
\begin{ltrdiagram}
% arch: TikZ / pgfplots source, built inside ltrdiagram (LTR)
\begin{tikzpicture}[font=\sffamily\small,
  bx/.style={draw=acc,fill=soft,rounded corners=4pt,minimum width=3.3cm,minimum height=1.2cm,align=center},
  hl/.style={draw=acc2,fill=soft2,rounded corners=4pt,minimum width=3.3cm,minimum height=1.2cm,align=center},
  ar/.style={-{Stealth[length=6pt]},thick,ink!70}]
\node[bx] (in) at (0,0) {\he{\textbf{קליטה}}\\\he{טביעה ונרמול}};
\node[bx] (l1) at (4.2,2.0) {\he{\textbf{שכבה 1: מקור}}\\\he{\en{C2PA} וסימני מים}};
\node[bx] (l2) at (4.2,0) {\he{\textbf{שכבה 2: זיהוי}}\\\he{4 גלאים בלתי תלויים}};
\node[bx] (l3) at (4.2,-2.0) {\he{\textbf{שכבה 3: הקשר}}\\\he{מקור ראשון, מטא דאטה}};
\node[hl] (fu) at (8.4,0) {\he{\textbf{שכבה 4: שקלול}}\\\he{ציון מכויל 0 עד 100}};
\node[bx] (ui) at (12.2,1.1) {\he{\textbf{מסוף העורך}}\\\he{החלטה ונימוק}};
\node[bx] (lg) at (12.2,-1.1) {\he{\textbf{יומן ודיווח}}\\\he{תיעוד וחתימה}};
\draw[ar] (in)--(l1); \draw[ar] (in)--(l2); \draw[ar] (in)--(l3);
\draw[ar] (l1)--(fu); \draw[ar] (l2)--(fu); \draw[ar] (l3)--(fu);
\draw[ar] (fu)--(ui); \draw[ar] (ui)--(lg);
\draw[ar,dashed,acc2] (l1.north) to[bend left=25] node[above,font=\sffamily\scriptsize]{\he{מסלול ירוק: תעודה תקפה}} (ui.north);
\end{tikzpicture}

\end{ltrdiagram}
\caption{שלוש שכבות מייצרות ראיות, שכבה אחת משקללת, ואדם מחליט}\label{fig:arch}
\end{figure}

טבלה~\ref{tab:subsys} מפרקת את השכבות שבאיור לתתי מערכות, ולכל אחת מגדירה תפקיד, קלט, פלט וטכנולוגיה.


\begin{table}[!htb]\centering
\caption{תתי המערכות: תפקיד, קלט, פלט וטכנולוגיה}\label{tab:subsys}
\small
\begin{tabularx}{\linewidth}{@{}l Y Y l@{}}
\toprule
\textbf{תת מערכת} & \textbf{תפקיד} & \textbf{פלט} & \textbf{טכנולוגיה} \\
\midrule
קליטה & טביעת \en{SHA-256}, נרמול פורמט, שמירת מקור & תיק תמונה & \en{Go, S3} \\
שכבה \num{1}: מקור & אימות תעודה, שרשרת עריכות, ביטולים & ממצא ״חתום / לא חתום / פגום״ & \en{c2pa-rs} \\
שכבה \num{2}: זיהוי & ארבעה גלאים: תחום תדר, רעש חיישן, עקביות תאורה, מסווג דיפוזיה & ארבעה ציונים ומפת חום & \en{PyTorch, GPU} \\
שכבה \num{3}: הקשר & חיפוש הפוך, השוואת מטא דאטה לכיתוב & הופעה ראשונה, סתירות & ספקים חיצוניים \\
שכבה \num{4}: שקלול & מודל כיול שמשלב את כל הממצאים & ציון \num{0} עד \num{100} והמלצה & \en{Isotonic Calibration} \\
מסוף העורך & הצגת ראיות והחלטה & החלטה חתומה & \en{React, RTL} \\
יומן ודיווח & רישום בלתי ניתן לשינוי, דוחות & יומן ודוח חודשי & \en{PostgreSQL, Merkle} \\
\bottomrule
\end{tabularx}
\end{table}

\section{זרימת ההחלטה}
איור~\ref{fig:flow} מגדיר את הנוהל המחייב. הוא מקודד במערכת ומקביל לכלל \num{4} בתקן \en{TRUTH-IMG}.

\begin{figure}[!htb]
\begin{ltrdiagram}
% flow: TikZ / pgfplots source, built inside ltrdiagram (LTR)
\begin{tikzpicture}[font=\sffamily\scriptsize,node distance=0.55cm and 0.9cm,
  st/.style={draw=acc,fill=soft,rounded corners=3pt,minimum width=2.7cm,minimum height=0.8cm,align=center},
  dc/.style={diamond,draw=acc2,fill=soft2,aspect=2.1,inner sep=1pt,align=center},
  ok/.style={st,draw=okc,fill=okc!12},
  md/.style={st,draw=midc!80!black,fill=midc!18},
  bd/.style={st,draw=badc,fill=badc!10},
  ar/.style={-{Stealth[length=5pt]},thick,ink!70}]
\node[st] (s) {\he{תמונה נקלטה}};
\node[dc,below=of s] (d1) {\he{תעודת \en{C2PA}}\\\he{תקפה?}};
\node[ok,right=1.6cm of d1] (g) {\he{מסלול ירוק:}\\\he{בדיקת הקשר בלבד}};
\node[st,below=of d1] (run) {\he{הרצת שכבות 2 עד 3}\\\he{וחישוב ציון}};
\node[dc,below=of run] (d2) {\he{ציון?}};
\node[ok,below left=0.7cm and 1.1cm of d2] (a) {\he{0 עד 29: ״מאומת״}\\\he{עורך מאשר}};
\node[md,below=0.9cm of d2] (b) {\he{30 עד 79: ״לא ודאי״}\\\he{בודק אנושי, עד 15 דק׳}};
\node[bd,below right=0.7cm and 1.1cm of d2] (c) {\he{80 עד 100: ״חשד גבוה״}\\\he{עורך בכיר בלבד}};
\node[st,below=1.9cm of b] (lab) {\he{תיוג לפי \en{TRUTH-IMG}}\\\he{ורישום ביומן}};
\draw[ar] (s)--(d1);
\draw[ar] (d1)--node[above]{\he{כן}}(g);
\draw[ar] (d1)--node[left]{\he{לא}}(run);
\draw[ar] (run)--(d2);
\draw[ar] (d2)-|(a); \draw[ar] (d2)--(b); \draw[ar] (d2)-|(c);
\draw[ar] (a)|-(lab); \draw[ar] (b)--(lab); \draw[ar] (c)|-(lab);
\draw[ar] (g)|-(lab);
\end{tikzpicture}

\end{ltrdiagram}
\caption{כל תמונה עוברת את אותה זרימה, ותמונה מפוקפקת מגיעה תמיד לאדם}\label{fig:flow}
\end{figure}

\section{תפיסת הפעלה: משמרת לילה}
כדי לבדוק שהזרימה עובדת במציאות ולא רק בתרשים, הנה תרחיש מלא, בזמנים שהמערכת מתחייבת אליהם:
\begin{enumerate}
\item \textbf{\num{23:41}}\quad שריפה גדולה בנמל. עורכת המשמרת מקבלת מקורא, בווטסאפ, תמונה דרמטית של להבות מעל מנופים. שלושה אתרים מתחרים כבר פרסמו תמונה דומה.
\item \textbf{\num{23:42}}\quad היא גוררת את התמונה למסוף. אין תעודת \en{C2PA}, כי המטא דאטה נמחק בשליחה.
\item \textbf{\num{23:42:50}}\quad שכבות \num{2} עד \num{3} מסתיימות. ציון \num{64}, ״לא ודאי״. מפת החום מסמנת את המנופים: הקווים שלהם ״נמסים״ זה לתוך זה. החיפוש ההפוך לא מוצא את התמונה בשום מקום לפני \num{23:30}, ואין מי שצילם אותה בשטח. אין סתירה מוכחת, אבל גם אין אישור.
\item \textbf{\num{23:44}}\quad לפי כלל \num{6}, התמונה אינה מתפרסמת כתיעוד. העורכת מפרסמת את הידיעה עם צילום ארכיון של הנמל, מתויג ״צילום ארכיון״, ומבקשת מהצלם בשטח צילום חתום.
\item \textbf{\num{23:58}}\quad מגיע צילום חתום מהצלם. מסלול ירוק, ציון \num{4}, ״מאומת״. הוא מחליף את צילום הארכיון.
\item \textbf{\num{00:20}}\quad הדסק מפרסם כתבת הפרכה קצרה על התמונה המזויפת שהופצה, לפי כלל \num{9}.
\end{enumerate}
התוצאה: הידיעה עלתה שלוש דקות אחרי המתחרים, והתמונה האמיתית עלתה \num{17} דקות אחריהם. בשעה \num{00:15} התברר שהתמונה מהקורא נוצרה בבינה מלאכותית, ומערכת החדשות הפכה את הזיוף עצמו לסיפור. זה הערך שהמוצר מוכר.

\section{מנוע השקלול והספים}
\label{sec:fusion}
כל גלאי מחזיר הסתברות גולמית. מנוע השקלול משלב אותן עם הממצאים משכבות \num{1} ו-\num{3} באמצעות רגרסיה לוגיסטית, ואחריה כיול איזוטוני על מאגר כיול נפרד:
\begin{equation*}
s = 100\cdot \mathrm{Cal}\!\left(\sigma\Big(\beta_0+\sum_{k=1}^{4}\beta_k\,\mathrm{logit}(p_k)+\gamma_1\,\mathbf{1}_{\text{wm}}+\gamma_2\,\mathbf{1}_{\text{ctx}}\Big)\right)
\end{equation*}
כאשר $p_k$ הוא פלט הגלאי ה-$k$, $\mathbf{1}_{\text{wm}}$ מציין שנמצא סימן מים של מחולל, ו-$\mathbf{1}_{\text{ctx}}$ מציין סתירה בהקשר (למשל, התמונה הופיעה ברשת לפני מועד האירוע הנטען). הכיול מבטיח שהציון ניתן לפירוש כהסתברות: מבין התמונות שקיבלו \num{80}, כ-\num{80\%} אכן מזויפות.

בהתפלגות המתוכננת, \num{91\%} מהתמונות הסינתטיות מקבלות \num{80} ומעלה (עומד ביעד \req{G-1}), ורק \num{1\%} מהאמיתיות. הספים \num{30} ו-\num{80} נבחרו כך שאזור ״לא ודאי״ יכיל לכל היותר \num{15\%} מהתמונות. כך העומס על הבודקים האנושיים נשאר סביר, וכמעט כל טעויות המודל מגיעות לאדם. איור~\ref{fig:dist} מציג את ההתפלגות הצפויה. זהו תכנון, ובפרויקט יוחלף בנתונים אמיתיים מהפיילוט.

\begin{figure}[!htb]
\begin{ltrdiagram}
% dist: TikZ / pgfplots source, built inside ltrdiagram (LTR)
\begin{tikzpicture}
\begin{axis}[width=12cm,height=5.2cm,ybar,bar width=7pt,
  xmin=-2,xmax=102,ymin=0,ymax=84,
  xlabel={\he{ציון \en{VeriLens}}}, ylabel={\he{\% מכל סוג}},
  xlabel style={font=\sffamily\small}, ylabel style={font=\sffamily\small},
  ticklabel style={font=\sffamily\scriptsize},
  legend style={font=\sffamily\scriptsize,at={(0.5,0.84)},anchor=north,legend columns=2,draw=none},
  area legend,
  axis x line*=bottom, axis y line*=left]
\fill[okc!8] (axis cs:-2,0) rectangle (axis cs:30,84);
\fill[midc!14] (axis cs:30,0) rectangle (axis cs:80,84);
\fill[badc!8] (axis cs:80,0) rectangle (axis cs:102,84);
\addplot[fill=acc!55,draw=acc] coordinates {(5,45)(15,30)(25,13)(35,5)(45,2.5)(55,1.5)(65,1)(75,1)(85,0.6)(95,0.4)};
\addplot[fill=acc2!75,draw=acc2] coordinates {(5,0.5)(15,0.5)(25,1)(35,1)(45,1)(55,1.5)(65,1.5)(75,2)(85,16)(95,75)};
\legend{\he{תמונות אמיתיות},\he{תמונות סינתטיות}}
\node[font=\sffamily\scriptsize] at (axis cs:14,79) {\he{״מאומת״}};
\node[font=\sffamily\scriptsize] at (axis cs:55,79) {\he{״לא ודאי״}};
\node[font=\sffamily\scriptsize] at (axis cs:91,79) {\he{״חשד גבוה״}};
\end{axis}
\end{tikzpicture}

\end{ltrdiagram}
\caption{היסטוגרמת ציונים מתוכננת: רוב התמונות מוכרעות אוטומטית, ופחות מ-\num{15\%} מגיעות לבודק}\label{fig:dist}
\end{figure}

\section{תווית לקורא}
המערכת מצמידה לכל תמונה אחת מארבע תוויות, לפי הסיווג בטבלה~\ref{tab:taxonomy}. לחיצה על התווית פותחת את תעודת המקור ואת רשומת ההחלטה, בלי פרטים מזהים של הבודק (איור~\ref{fig:labels}).

\begin{figure}[!htb]
\begin{ltrdiagram}
% labels: TikZ / pgfplots source, built inside ltrdiagram (LTR)
\begin{tikzpicture}[font=\sffamily\small,
  lb/.style={rounded corners=9pt,minimum height=0.75cm,inner sep=6pt,text=white,font=\sffamily\small\bfseries}]
\node[lb,fill=okc] (a) at (0,0) {\he{צילום מאומת}};
\node[lb,fill=acc] (b) at (3.9,0) {\he{צילום ערוך}};
\node[lb,fill=midc!85!black] (c) at (7.6,0) {\he{איור סינתטי}};
\node[lb,fill=badc] (d) at (11.2,0) {\he{תמונה שזוהתה כמזויפת}};
\node[font=\sffamily\scriptsize,align=center] at (0,-0.85) {\he{קטגוריה א׳}};
\node[font=\sffamily\scriptsize,align=center] at (3.9,-0.85) {\he{קטגוריה ב׳, עריכה מוצהרת}};
\node[font=\sffamily\scriptsize,align=center] at (7.6,-0.85) {\he{קטגוריה ד׳}};
\node[font=\sffamily\scriptsize,align=center] at (11.2,-0.85) {\he{רק בכתבה שמפריכה את הזיוף}};
\end{tikzpicture}

\end{ltrdiagram}
\caption{ארבע תוויות לקורא: הקורא יודע תמיד מה הוא רואה}\label{fig:labels}
\end{figure}

\section{התקן האתי \en{TRUTH-IMG}}\label{sec:truthimg}
התקן מתרגם את העקרונות של \en{SPJ}, \en{AP} ואמנת פריז~\cite{spj2014,ap2023,paris2023} לעשרה כללים מחייבים. כל כלל ניתן לביקורת, כלומר אפשר לבדוק בדיעבד אם קוים (טבלה~\ref{tab:truth}).

\begin{table}[!htb]\centering
\caption{עשרת הכללים של \en{TRUTH-IMG} גרסה \num{1.0}}\label{tab:truth}
\small
\begin{tabularx}{\linewidth}{@{}c Y Y@{}}
\toprule
\textbf{\#} & \textbf{כלל} & \textbf{איך מבקרים} \\
\midrule
\num{1} & תמונה סינתטית (קטגוריה ג׳) לעולם לא תוצג כתיעוד. & סריקת ארכיון מדגמית \\
\num{2} & אין לשנות צילום חדשותי בכלים גנרטיביים; עריכה טכנית בלבד. & השוואת קובץ מקור לקובץ שפורסם \\
\num{3} & איור בינה מלאכותית (קטגוריה ד׳) מתויג בתווית גלויה ובכיתוב. & סריקה אוטומטית של תוויות \\
\num{4} & כל תמונה עוברת את זרימת ההחלטה (איור~\ref{fig:flow}) לפני פרסום. & השוואת יומן לתמונות שפורסמו \\
\num{5} & ההחלטה הסופית היא תמיד של אדם מזוהה. & יומן: שדה ״מאשר״ מלא \\
\num{6} & ״לא ודאי״ פירושו ״לא מפרסמים כתיעוד״ עד שהספק נפתר. & יומן: אין פרסום בלי אישור בכיר \\
\num{7} & מהירות או טראפיק אינם נימוק לגיטימי לפרסום תמונה לא מאומתת. & סקירת נימוקים ביומן \\
\num{8} & טעות שהתגלתה תתוקן בגלוי תוך \num{24} שעות, עם הסבר לקורא. & דוח תיקונים חודשי \\
\num{9} & בכתבה שמפריכה זיוף, התמונה המזויפת מוצגת מסומנת ומוקטנת, ולא כתמונה ראשית. & בחינת כתבות הפרכה \\
\num{10} & דוח שקיפות ציבורי חודשי: מספר התמונות, החלטות, תיקונים. & פרסום הדוח באתר \\
\bottomrule
\end{tabularx}
\end{table}

כלל \num{7} הוא לב התקן: הוא מנתק במפורש את הקשר בין התמריץ הכלכלי לבין ההחלטה העריכתית. זהו המקום שבו האתיקה העיתונאית פוגשת את מודל הטראפיק. מבחינה דאונטולוגית, פרסום תמונה מפוקפקת כדי למשוך קליקים מתייחס לקורא כאמצעי להכנסה ולא כמי שזכאי לאמת. מבחינה תוצאתית, הרווח הקצר בטראפיק קטן מהנזק המצטבר לאמון.

\section{ארכיטקטורה לקנה מידה}
המעבר מ-\num{2} מיליון ל-\num{50} מיליון תמונות בחודש (פי \num{25}) אינו דורש שינוי בארכיטקטורה, אלא הוספת משאבים. כל שכבה היא שירות נפרד מעל תור הודעות, והגלאים רצים על מאגר \en{GPU} שגדל אוטומטית לפי עומס. התמונות מנותבות לאזור אחסון לפי מיקום הלקוח (אירופה, ישראל, ארה״ב), בהתאם ל-\req{SR-04} (טבלה~\ref{tab:capacity}).

\begin{table}[H]\centering
\caption{תכנון קיבולת: מבחן עומס מול יעד \num{36} חודשים}\label{tab:capacity}
\begin{tabularx}{\linewidth}{@{}Y c c@{}}
\toprule
\textbf{פרמטר} & \textbf{מבחן עומס (חודש \num{16})} & \textbf{יעד (חודש \num{36})} \\
\midrule
תמונות בחודש & \num{2} מיליון & \num{50} מיליון \\
שיא בשעה (פי \num{4} מהממוצע) & \num{11,000} & \num{280,000} \\
מעבדי \en{GPU} פעילים בשיא & \num{6} & \num{120} \\
עלות מחשוב לתמונה (יעד) & \num{0.04} ₪ & \num{0.012} ₪ \\
אזורי אחסון & \num{1} & \num{3} \\
\bottomrule
\end{tabularx}
\end{table}
